\section{平坦性}\label{src:A.2}
\begin{epigraph}
艺术是火焰加上代数.
\epigraphauthor{Jorge Luis Borges}
\end{epigraph}
没有平坦性, 模理论就寸步难行. 尽管其定义看似抽象且纯属代数, 它却是态射 $X\to Y$ 的一种神奇的几何性质, 保证当 $y\in Y$ 变化时, 纤维 $X_y$ ``良好地变化''. Hilbert 多项式对平坦性的刻画 (\bookref{A.2.4}) 很好地体现了这一原则. 这正是我们把模叠的对象定义为平坦族的原因.

\subsection{平坦性判别法}\label{src:A.2.1-subsection}
称环 $A$ 上的模 $M$ \newterm{平坦}{flat}, 如果函子
\[
-\otimes_A M:\Mod(A)\longrightarrow\Mod(A)
\]
正合. 我们回顾以下判别法:
\begin{enumerate}[(1)]
\item (茎判别法) $M$ 在 $A$ 上平坦, 当且仅当对每个素理想 (等价地, 极大理想) $\mathfrak p$, $M_{\mathfrak p}$ 在 $A_{\mathfrak p}$ 上平坦. 更一般地, 若 $A\to B$ 为环映射, 则 $B$-模 $N$ 平坦, 当且仅当对每个原像为 $\mathfrak p\subseteq A$ 的素理想 $\mathfrak q\subseteq B$, $N_{\mathfrak q}$ 在 $A_{\mathfrak p}$ 上平坦.
\item (理想判别法) $M$ 平坦, 当且仅当对每个有限生成理想 $I\subseteq A$, 映射 $I\otimes_A M\to M$ 为单射 \cite[命题 $6.1$]{Eis95}. (当 $A$ 为主理想整环时, 这意味着 $M$ 平坦当且仅当 $M$ 无挠.)
\item ($\operatorname{Tor}$ 判别法) $M$ 平坦, 当且仅当对所有有限生成理想 $I\subseteq A$, $\operatorname{Tor}_1^A(A/I,M)=0$ \cite[命题 $6.1$]{Eis95}.
\item (有限表现判别法) $M$ 在 $A$ 上有限表现且平坦, 当且仅当 $M$ 有限且投射, 也当且仅当 $M$ 有限且局部自由 (即对所有素理想, 等价地所有极大理想 $\mathfrak p$, $M_{\mathfrak p}$ 有限且自由); 见 \cite[Tag 00NX]{SP}. (没有有限表现假设时, Lazard 定理断言, $M$ 在 $A$ 上平坦, 当且仅当 $M$ 可以写成有限自由 $A$-模 $M_i$ 的有向余极限 $\operatorname{colim}_{i\in I}M_i$; 见 \cite[A$6.6$]{Eis95} 或 \cite[Tag 058G]{SP}.)
\item (方程判别法) $M$ 平坦, 当且仅当对每个关系 $\sum_{i=1}^n a_i m_i=0$, 其中 $a_i\in A$ 和 $m_i\in M$, 存在 $m'_j\in M$ ($j=1,\ldots,r$) 和 $a'_{ij}\in A$, 使得对所有 $i$ 有 $\sum_{j=1}^r a'_{ij}m'_j=m_i$, 对所有 $j$ 有 $\sum_{i=1}^n a'_{ij}a_i=0$ \cite[推论 $6.5$]{Eis95}.
\end{enumerate}
若 $f:X\to Y$ 为概形态射, 则称拟凝聚 $\mathcal O_X$-模 $\mathcal F$ 平坦, 如果对所有仿射开集 $\Spec B\subseteq Y$ 和 $\Spec A\subseteq f^{-1}(\Spec B)$, $B$-模 $\Gamma(\Spec A,\mathcal F)$ 平坦.
\begin{proposition}[平坦性的等价条件]\label{src:A.2.1}
设 $f:X\to Y$ 为概形态射, $\mathcal F$ 为拟凝聚 $\mathcal O_X$-模. 下列条件等价:
\begin{enumerate}[(1)]
\item $\mathcal F$ 在 $Y$ 上平坦;
\item 存在 $Y$ 的 Zariski 覆盖 $\{\Spec B_i\}$ 和 $f^{-1}(\Spec B_i)$ 的 Zariski 覆盖 $\{\Spec A_{ij}\}$, 使得 $\Gamma(\Spec A_{ij},\mathcal F)$ 在环映射 $B_i\to A_{ij}$ 下为平坦 $B_i$-模;
\item 对所有 $x\in X$, $\mathcal O_{X,x}$-模 $\mathcal F_x$ 作为 $\mathcal O_{Y,y}$-模平坦;
\item 函子
\[
\QCoh(Y)\longrightarrow\QCoh(X),\qquad\mathcal G\longmapsto f^*\mathcal G\otimes_{\mathcal O_X}\mathcal F
\]
正合.
\end{enumerate}
\end{proposition}
\begin{proof}
见 \cite[\S III.$9$]{Har77} 或 \cite[Tag 01U2]{SP}.
\end{proof}
称概形态射 $f:X\to Y$ 在 $x\in X$ 处\textit{平坦} (相应地, 拟凝聚 $\mathcal O_X$-模 $\mathcal F$ 在 $x\in X$ 处\textit{平坦}), 如果存在包含 $x$ 的 Zariski 开邻域 $U\subseteq X$, 使得 $f|_U$ (相应地, $\mathcal F|_U$) 在 $Y$ 上平坦. 这等价于 $\mathcal O_{X,x}$ (相应地, $\mathcal F_x$) 作为 $\mathcal O_{Y,y}$-模平坦.
\begin{proposition}[光滑曲线上的平坦性判别法]\label{src:A.2.2}
设 $C$ 为维数为 $1$ 的整正则概形 (例如离散赋值环的谱, 或域上的光滑连通曲线), $X\to C$ 为概形态射. 拟凝聚 $\mathcal O_X$-模 $\mathcal F$ 在 $C$ 上平坦, 当且仅当 $\mathcal F$ 的每个相伴点都映到 $C$ 的一般点.
\end{proposition}
\begin{proof}
由离散赋值环上的模平坦当且仅当其无挠, 经一个简短论证即可得到结论; 见 \cite[III.$9.7$]{Har77}.
\end{proof}
在高维基底上, 有时可以归结为上述光滑曲线上的判别法来检验平坦性. 这称为\newterm{平坦性的赋值判别法}{valuative criterion for flatness}: 若 $f:X\to S$ 为 Noether 概形之间的有限型态射, $S$ 约化, $\mathcal F$ 为凝聚 $\mathcal O_X$-模, 则 $\mathcal F$ 在 $x\in X$ 处平坦, 当且仅当对每个来自离散赋值环的映射 $(\Spec R,0)\to(S,f(x))$, 限制 $\mathcal F|_{X_R}$ 在 $X_R:=X\times_S\Spec R$ 中位于 $0$ 和 $x$ 上方的所有点处都在 $R$ 上平坦 \cite[IV.$11.8.1$]{EGA}. 尽管它给出了一个概念性的几何平坦性判别法, 令人意外的是, 它在模理论中很少使用.
\begin{proposition}[Artin 环上的平坦性判别法]\label{src:A.2.3}
Artin 环上的模平坦, 当且仅当它自由, 也当且仅当它投射.
\end{proposition}
\begin{proof}
见 \cite[Tag 051E]{SP}.
\end{proof}
回顾, 若 $X\subseteq\mathbb P_K^n$ 为子概形, $\mathcal F$ 为拟凝聚 $\mathcal O_X$-模, 则 $\mathcal F$ 的 Hilbert 多项式为 $P_{\mathcal F}(z)=\chi(X,\mathcal F(z))\in\mathbb Q[z]$.
\begin{proposition}[通过 Hilbert 多项式刻画平坦性]\label{src:A.2.4}
设 $S$ 为约化, 连通的 Noether 概形, $X\subseteq\mathbb P_S^n$ 为闭子概形. 凝聚 $\mathcal O_X$-模 $\mathcal F$ 在 $S$ 上平坦, 当且仅当函数
\[
S\longrightarrow\mathbb Q[z],\qquad s\longmapsto P_{\mathcal F|_{X_s}}
\]
为常值; 这里它把点 $s\in S$ 映到限制 $\mathcal F|_{X_s}$ 在纤维 $X_s\subseteq\mathbb P_{\kappa(s)}^n$ 上的 Hilbert 多项式.
\end{proposition}
\begin{proof}
见 \cite[定理 $9.9$]{Har77}.
\end{proof}
\begin{theorem}[平坦性的局部与无穷小判别法]\label{src:A.2.5}
设 $A\to B$ 为 Noether 局部环之间的局部同态, $M$ 为有限 $B$-模. 下列条件等价:
\begin{enumerate}[(1)]
\item $M$ 在 $A$ 上平坦;
\item (局部判别法) $\operatorname{Tor}_1^A(A/\mathfrak m_A,M)=0$;
\item (无穷小判别法) 对每个 $n\geq1$, $M/\mathfrak m_A^n M$ 在 $A/\mathfrak m_A^n$ 上平坦.
\end{enumerate}
\end{theorem}
\begin{proof}
见 \cite[定理 $6.8$, 习题 $6.5$]{Eis95} 或 \cite[Tag 00MK]{SP}.
\end{proof}
平坦性的局部判别法的下述推论在形变理论中特别有用.
\begin{corollary}\label{src:A.2.6}
设 $A\to A_0$ 为 Noether 环之间的满同态, 其核 $I$ 满足 $I^2=0$. $A$-模 $M$ 在 $A$ 上平坦, 当且仅当
\begin{enumerate}[(1)]
\item $M_0:=M\otimes_A A_0$ 在 $A_0$ 上平坦;
\item 映射 $M_0\otimes_{A_0}I\to M$ 为单射.
\end{enumerate}
\end{corollary}
\begin{proof}
对于 ($\Rightarrow$), 条件 (1) 由基变换成立; 把正合列 $0\to I\to A\to A_0\to0$ 与 $M$ 作张量积, 并使用等同 $M\otimes_A I\simeq M_0\otimes_{A_0}I$, 即得条件 (2). 对于 ($\Leftarrow$), 由平坦性的局部判别法 (\bookref{A.2.5}), 只需证明对所有素理想 $\mathfrak p\subseteq A$, $\operatorname{Tor}_1^A(A/\mathfrak p,M)=0$. 令 $\mathfrak p_0:=\mathfrak p/I\subseteq A$. 考虑把正合列 $0\to I\to\mathfrak p\to\mathfrak p_0\to0$ 和 $0\to I\to A\to A_0\to0$ 与 $M$ 作张量积得到的图:
\[
\begin{tikzcd}[column sep=1.1em,row sep=1.8em]
&&\operatorname{Tor}_1^A(M,A/\mathfrak p)\arrow[r]\arrow[d]&\operatorname{Tor}_1^{A_0}(M_0,A_0/\mathfrak p_0)\arrow[d]&&\\
0\arrow[r]&M\otimes_A I\arrow[r]\arrow[d,equal]&M\otimes_A\mathfrak p\arrow[r]\arrow[d]&M_0\otimes_{A_0}\mathfrak p_0\arrow[r]\arrow[d]&0\\
0\arrow[r]&M\otimes_A I\arrow[r]&M\arrow[r]\arrow[d,two heads]&M_0\arrow[r]\arrow[d,two heads]&0\\
&&M\otimes_A A/\mathfrak p\arrow[r]&M_0\otimes_{A_0}A_0/\mathfrak p_0
\end{tikzcd}
\]
条件 (2) 意味着第二行正合, 进而第一行也正合; 这里使用了等同 $M\otimes_A\mathfrak p_0\simeq M_0\otimes_{A_0}\mathfrak p_0$. 条件 (1) 意味着 $\operatorname{Tor}_1^{A_0}(M_0,A_0/\mathfrak p_0)=0$, 再由蛇引理得 $\operatorname{Tor}_1^A(M,A/\mathfrak p)=0$. 另见 \cite[命题 $2.2$]{Har10}.
\end{proof}
\begin{remark}\label{src:A.2.7}
对双数环 $A=\Bbbk[\epsilon]/(\epsilon^2)$ 和 $A'=\Bbbk$ 应用上述结果, 得到 $A$-模 $M$ 平坦当且仅当 $M\otimes_{\Bbbk[\epsilon]/(\epsilon^2)}\Bbbk\xrightarrow{\epsilon}M$ 为单射. 这也可以由以下事实推出: 环 $B$ 上的模 $N$ 平坦, 当且仅当对每个理想 $I\subseteq B$, 映射 $I\otimes_B M\to M$ 为单射; 并注意 $\Bbbk[\epsilon]/(\epsilon^2)$ 中唯一的理想为 $(\epsilon)$.
\end{remark}
下面这些方便的事实与平坦性的局部判别法 (\bookref{A.2.5}) 密切相关.
\begin{lemma}\label{src:A.2.8}
设 $(A,\mathfrak m_A)\to(B,\mathfrak m_B)$ 为 Noether 局部环之间的局部环同态.
\begin{enumerate}[(1)]
\item 设 $M$ 为平坦 $A$-模, $N$ 为有限生成 $B$-模. 若 $\phi:N\to M$ 为 $R$-模态射, 且 $N/\mathfrak m N\to M/\mathfrak m M$ 为单射, 则 $\phi:N\to M$ 为单射, 且 $M/\phi(N)$ 在 $A$ 上平坦.
\item 若 $A\to B$ 还平坦, 且 $f\in\mathfrak m_B$ 在 $B\otimes_A A/\mathfrak m_A$ 中为非零因子, 则 $A\to B/(f)$ 平坦.
\end{enumerate}
\end{lemma}
\begin{proof}
证明是初等的 (见 \cite[定理 $22.5$]{Mat89} 或 \cite[Tag 00ME]{SP}). 注意, 第 (2) 条直接由第 (1) 条推出.
\end{proof}
第 (2) 条可以视为 ``平坦性的切片判别法'', 并常对正则序列归纳应用. 它有如下几何解释.
\begin{corollary}[平坦性的切片判别法]\label{src:A.2.9}
设 $f:X\to S$ 为局部有限表现态射, $x\in X$ 的像为 $s\in S$. 若 $f$ 在 $x$ 处平坦, 且 $h\in\mathfrak m_x\subseteq\mathcal O_{X,x}$ 在纤维的局部环 $\mathcal O_{X_s,x}$ 中的像为非零因子, 则存在 $x$ 的开邻域 $U\subseteq X$, 使得 $h$ 延伸为 $U$ 上的全局函数, 且复合 $V(h)\hookrightarrow U\to S$ 局部有限表现并在 $x$ 处平坦.
\end{corollary}
\begin{proof}
Noether 情形是引理 \bookref{A.2.8}(2) 的直接推论, 一般情形可用 \S\bookref{B.3} 的极限方法归结为 Noether 情形. 另见 \cite[Tag 056X]{SP}.
\end{proof}
\begin{theorem}[纤维平坦性判别法]\label{src:A.2.10}
考虑概形的交换图
\[
\begin{tikzcd}[column sep=2em,row sep=1.5em]
X\arrow[rr]\arrow[dr]&&Y\arrow[dl]\\&S&
\end{tikzcd}
\]
设 $F$ 为有限表现拟凝聚 $\mathcal O_X$-模. 假设 $X\to S$ 局部有限表现, $Y\to S$ 局部有限型. 设 $x\in X$ 的像为 $y\in Y$ 和 $s\in S$. 若茎 $\mathcal F_x$ 非零, 则下列条件等价:
\begin{enumerate}[(1)]
\item $F$ 在 $x$ 处在 $S$ 上平坦, 且 $\mathcal F_s:=\mathcal F|_{X_s}$ 在 $x$ 处在 $Y_s$ 上平坦;
\item $Y$ 在 $y$ 处在 $S$ 上平坦, 且 $F$ 在 $x$ 处在 $Y$ 上平坦.
\end{enumerate}
\end{theorem}
\begin{proof}
见 \cite[Tag 039A]{SP}.
\end{proof}
若 $A\to B$ 为 Noether 局部环之间的局部环映射, 则 $\dim B=\dim A+\dim B/\mathfrak m_A B$. 以下是一个部分逆命题.
\begin{theorem}[奇迹平坦性]\label{src:A.2.11}
设 $A\to B$ 为 Noether 局部环之间的局部同态. 假设
\begin{enumerate}[1.]
\item $A$ 正则;
\item $B$ 为 Cohen--Macaulay 环;
\item $\dim B=\dim A+\dim B/\mathfrak m_A B$.
\end{enumerate}
则 $A\to B$ 平坦.
\end{theorem}
\begin{proof}
见 \cite[定理 $25.16$]{Nag62} 或 \cite[Tag 00R4]{SP}.
\end{proof}

\subsection{平坦性的性质}\label{src:A.2.2-subsection}
\begin{proposition}[平坦基变换]\label{src:A.2.12}
考虑概形的 Cartesian 图
\[
\begin{tikzcd}[column sep=2.5em,row sep=2em]
X'\arrow[r,"g'"]\arrow[d,"f'"']\arrow[dr,phantom,"\square" description]&X\arrow[d,"f"]\\
Y'\arrow[r,"g"']&Y
\end{tikzcd}
\]
设 $F$ 为 $X$ 上的拟凝聚层. 若 $g:Y'\to Y$ 平坦, $f:X\to Y$ 拟紧且拟分离, 则自然伴随映射
\[
g^*\mathrm R^i f_*F\longrightarrow\mathrm R^i f'_*g'^*F
\]
对所有 $i\geq0$ 都是同构.
\end{proposition}
\begin{proof}
见 \cite[命题 III.$9.3$]{Har77} 或 \cite[Tag 02KH]{SP}.
\end{proof}
\begin{theorem}[一般平坦性]\label{src:A.2.13}
设 $f:X\to S$ 为概形之间的有限型态射, $\mathcal F$ 为有限型拟凝聚 $\mathcal O_X$-模. 若 $S$ 约化, 则存在稠密开子概形 $U\subseteq S$, 使得 $X_U\to U$ 平坦且为表现态射, 并使 $\mathcal F|_{X_U}$ 在 $U$ 上平坦且作为 $\mathcal O_{X_U}$-模有限表现.
\end{theorem}
\begin{proof}
见 \cite[Tag 052B]{SP}.
\end{proof}
\begin{proposition}[fppf 态射是开映射]\label{src:A.2.14}
设 $f:X\to Y$ 为概形态射. 若 $f$ 平坦且局部有限表现, 则对每个开子集 $U\subseteq X$, 像 $f(U)\subseteq Y$ 为开集.
\end{proposition}
\begin{proof}
见 \cite[Tag 01UA]{SP}.
\end{proof}
\begin{proposition}\label{src:A.2.15}
局部有限表现的平坦单态射 (例如平展单态射) 是开浸入.
\end{proposition}
\begin{proof}
这是下降理论的一个漂亮应用. 由上一个结果, 只需假设 $f:X\to Y$ 满射, 因而是 fppf 态射. 由于 $X\to Y$ 为单态射, 基变换 $X\times_Y X\to Y$ 为同构. 由于同构是靶上的 fpqc 局部性质 (\bookref{3.1.26}), $X\to Y$ 为同构.
\end{proof}
下面的定理是 Quot 概形的射影性 (\bookref{2.1.3}) 的一个关键要素.
\begin{theorem}[平坦化分层的存在性]\label{src:A.2.16}
设 $X\to S$ 为 Noether 概形之间的射影态射, $\mathcal O_X(1)$ 为相对丰沛线丛, $\mathcal F$ 为 $X$ 上的凝聚层. 对每个多项式 $P\in\mathbb Q[z]$, 存在局部闭子概形 $S_P\hookrightarrow S$, 使得态射 $T\to S$ 经由 $S_P$ 分解, 当且仅当 $\mathcal F$ 到 $X_T$ 的拉回 $\mathcal F_T$ 在 $T$ 上平坦, 且对每个 $t\in T$, 到 $X_{\kappa(t)}$ 的拉回 $\mathcal F_{\kappa(t)}$ 的 Hilbert 多项式为 $P$.

此外, 存在由多项式组成的有限指标集 $I$, 使得在集合意义下 $S=\coprod_{P\in I}S_P$. $S_P$ 在 $S$ 中的闭包在集合意义下包含于并集 $\bigcup_{P\leq Q}S_Q$, 其中 $P\leq Q$ 当且仅当对 $z\gg0$ 有 $P(z)\leq Q(z)$.
\end{theorem}
\begin{proof}
见 \cite[引理 $3.3$]{FGAIV} 或 \cite[\S$8$]{Mum66a}. 核心思想是归结为 $X=S$, 此时闭子集 $\overline{S_P}$ 由 $\mathcal F$ 的 Fitting 理想定义 \cite[Tag 05P9]{SP}.
\end{proof}
\begin{remark}\label{src:A.2.17}
若 $X\to S$ 仅固有, 则存在\newterm{万有平坦化}{universal flattening}, 即存在代数空间 $S'$ 及态射 $S'\to S$, 使得映射 $T\to S$ 经由 $S'\to S$ 分解, 当且仅当 $\mathcal F$ 到 $X_T:=X\times_S T$ 的拉回 $\mathcal F|_{X_T}$ 在 $T$ 上平坦 \cite[Tag 05UG]{SP}. 一般而言, $S'$ 未必是 $S$ 的局部闭子概形的不交并; 见 \cite{Kre13}.
\end{remark}
\begin{theorem}[Raynaud--Gruson 平坦化]\label{src:A.2.18}
设 $Y$ 为拟紧且拟分离概形, $X\to Y$ 为有限表现态射, 并在拟紧开子概形 $U\subseteq Y$ 上平坦. 则存在交换图
\[
\begin{tikzcd}[column sep=2.5em,row sep=2em]
\widetilde X\arrow[r]\arrow[d]&X\arrow[d,"f"]\\
Y'\arrow[r,"p"']&Y
\end{tikzcd}
\]
其中 $p:Y'\to Y$ 为沿一个与 $U$ 不交的有限表现闭子概形 $Z\subseteq Y$ 的爆破, 且 $X$ 的严格变换 $\widetilde X$ 在 $Y'$ 上平坦.
\end{theorem}
上述严格变换 $\widetilde X$ 按定义是 $(Y'\setminus p^{-1}(Z))\times_Y X$ 在基变换 $Y'\times_Y X$ 中的闭包.
\begin{proof}
见 \cite[定理 I.$5.2.2$]{GR71} 或 \cite[Tag 0815]{SP}.
\end{proof}

\subsection{忠实平坦性}\label{src:A.2.3-subsection}
称环 $A$ 上的模 $M$ \newterm{忠实平坦}{faithfully flat}, 如果函子 $-\otimes_A M:\Mod(A)\to\Mod(A)$ 忠实正合, 即 $A$-模序列 $N'\to N\to N''$ 正合当且仅当 $N'\otimes_A M\to N\otimes_A M\to N''\otimes_A M$ 正合.

对每个非零 $A$-模映射 $\phi:N\to N'$, 诱导的映射 $\phi\otimes_A M:N\otimes_A M\to N'\otimes_A M$ 也非零.
\begin{proposition}[忠实平坦性的等价条件]\label{src:A.2.19}
设 $A$ 为环, $M$ 为平坦 $A$-模. 下列条件等价:
\begin{enumerate}[(1)]
\item $M$ 忠实平坦;
\item 对每个非零 $A$-模映射 $\phi:N\to N'$, 诱导映射 $\phi\otimes_A M:N\otimes_A M\to N'\otimes_A M$ 也非零;
\item 对每个非零 $A$-模 $N$, 张量积 $N\otimes_A M$ 非零;
\item 对每个素理想 $\mathfrak p\subseteq A$, 张量积 $M\otimes_A\kappa(\mathfrak p)$ 非零;
\item 对每个极大理想 $\mathfrak m\subseteq A$, 张量积 $M\otimes_A\kappa(\mathfrak m)\simeq M/\mathfrak m M$ 非零.
\end{enumerate}
\end{proposition}
\begin{proof}
见 \cite[Tag 00H9]{SP}.
\end{proof}
当 $M=B$ 为 $A$-代数时, 由 (4), 平坦环映射 $A\to B$ 忠实平坦, 当且仅当 $\Spec B\to\Spec A$ 满射; 或等价地, 由 (5), $A$ 的每个极大理想都在 $\Spec B\to\Spec A$ 的像中. 后一种等价性意味着任何平坦局部环映射都忠实平坦.

称概形态射 $f:X\to Y$ \textit{忠实平坦}, 如果 $f$ 平坦且满射. 这等价于 $f^*:\QCoh(Y)\to\QCoh(X)$ 忠实正合, 也等价于拟凝聚 $\mathcal O_Y$-模 (相应地, 拟凝聚 $\mathcal O_Y$-模的态射) 为零当且仅当其拉回为零. 忠实平坦态射在下降理论中起重要作用; 见 \S\bookref{3.1}.

\subsection{fppf 与 fpqc 态射}\label{src:A.2.4-subsection}
fppf 与 fpqc 分别是 ``fidèlement plate de présentation finie'' 和 ``fidèlement plat et quasi-compact'' 的缩写. 尽管如此命名, fpqc 态射比忠实平坦且拟紧的映射更一般.
\begin{definition}\label{src:A.2.20}
称概形态射 $f:X\to Y$ 为:
\begin{enumerate}[(1)]
\item \newterm{fppf 态射}{fppf}, 如果 $f$ 忠实平坦且局部有限表现;
\item \newterm{fpqc 态射}{fpqc}, 如果 $f$ 忠实平坦, 且 $Y$ 的每个拟紧开子集都是 $X$ 的某个拟紧开子集的像.
\end{enumerate}
\end{definition}
\begin{remark}\label{src:A.2.21}
拟紧且忠实平坦的态射是 fpqc 态射. 此外, 开且忠实平坦的态射也是 fpqc 态射: 对拟紧开子集 $V\subseteq Y$, 可以把 $f^{-1}(V)=\bigcup_i U_i$ 写成仿射开集的并; 由于每个 $f(U_i)\subseteq V$ 都开且 $V$ 拟紧, $V$ 是有限多个 $U_i$ 的像. fppf 态射是开映射 (\bookref{A.2.14}) 给出蕴含 $\mathrm{fppf}\Rightarrow\mathrm{fpqc}$.

fpqc 态射 $f:X\to Y$ 还可以等价地通过以下任一要求刻画: 存在 $Y$ 的仿射覆盖 $\{Y_i\}$, 使每个 $Y_i$ 都是 $X$ 的某个拟紧开子集的像; 或要求每个点 $x\in X$ 都有开 (相应地, 拟紧开) 邻域 $U$, 使得 $f(U)$ 开且 $U\to f(U)$ 拟紧; 见 \cite[命题 $2.33$]{Nit05}.

fppf (相应地, fpqc) 覆盖 $\{X_i\to X\}$ 是一族态射, 使得 $\coprod_i X_i\to X$ 为 fppf (相应地, fpqc) 态射.
\end{remark}

\subsection{泛单射同态}\label{src:A.2.5-subsection}
平坦模的定义性特征是, 与它作张量积保持所有单射. 对偶的概念是与任意模作张量积后仍保持为单射的模单射; 这也是一个很有用的性质.
\begin{definition}\label{src:A.2.22}
称 $A$-模同态 $M\to N$ 为\newterm{泛单射}{universally injective}, 如果对每个 $A$-模 $P$, 映射 $M\otimes_A P\to N\otimes_A P$ 为单射. 称环同态 $A\to B$ 泛单射, 如果它作为 $A$-模映射泛单射.
\end{definition}
\begin{remark}\label{src:A.2.23}
这不应与概形的泛单射或根式态射 $X\to Y$ 混淆, 后者指在任意基变换后仍为单射的单射映射; 见 \cite[Tag 01S2]{SP}.
\end{remark}
这一概念将在凝聚 Tannaka 对偶性 (定理 \bookref{7.6.1}) 的证明中发挥根本作用. 为此, 我们将使用以下性质:
\begin{proposition}\label{src:A.2.24}
\begin{enumerate}[(1)]
\item 忠实平坦环同态 $A\to B$ 泛单射.
\item $A$-模的分裂单射 $M\to N$ 泛单射. 若 $N/M$ 有限表现, 则逆命题成立.
\item 若 $A\to A'$ 忠实平坦, 则 $A$-模映射 $M\to N$ 泛单射, 当且仅当 $M\otimes_A A'\to N\otimes_A A'$ 泛单射.
\item 若 $A\to B$ 泛单射, 且由 $b\mapsto b\otimes1$ 定义的 $B\to B\otimes_A B$ 忠实平坦, 则 $A\to B$ 忠实平坦.
\end{enumerate}
\end{proposition}
\begin{proof}
第 (1), (2), (4) 条见 \cite[Tags 08WP, 058L, 08XD]{SP}. 第 (3) 条直接由 $-\otimes_A A'$ 的忠实正合性推出. 另见 \cite{Laz69} 或 \cite[\S$4$J]{Lam99}.
\end{proof}
值得注意的是, 泛单射环映射恰好是满足模的有效下降的映射; 见注 \bookref{3.1.6}.
